\section{¿Por qué el aprendizaje profundo? ¿Por qué ahora?}

\subsection{Las limitaciones del aprendizaje automático tradicional}

El aprendizaje automático tradicional depende de \textbf{características diseñadas a mano} (hand-engineered features). Tomando como ejemplo la detección de rostros, el ingeniero tiene que definir manualmente cómo extraer de la imagen características como bordes, curvas y rasgos faciales. Este método presenta tres problemas fundamentales:

\begin{enumerate}
\item \textbf{Consume tiempo}: la ingeniería de características requiere una gran cantidad de tiempo y experiencia de especialistas del dominio
\item \textbf{Es frágil}: las características diseñadas a mano no son robustas ante cambios del entorno (iluminación, ángulo, etc.)
\item \textbf{No es escalable}: cada tarea nueva exige diseñar de nuevo el conjunto de características
\end{enumerate}

\begin{importantbox}{La idea central del aprendizaje profundo}
El avance clave del aprendizaje profundo consiste en: \textbf{dejar que el modelo aprenda automáticamente, a partir de los datos, representaciones jerárquicas de las características} (hierarchical feature representations). Las características de bajo nivel (como bordes y líneas) se combinan automáticamente en características de nivel medio (como ojos, nariz), y estas a su vez se combinan en características de alto nivel (como la estructura facial). Sin intervención humana.
\end{importantbox}

\begin{figure}[H]
\centering
\includegraphics[width=0.85\textwidth]{slides-images/slide-018.jpg}
\caption{Aprendizaje jerárquico de características: de características de bajo nivel a características de alto nivel\protect\footnotemark}
\end{figure}
\footnotetext{Diapositivas, página 18.}

\subsection{Las tres fuerzas impulsoras}

Aunque la teoría de las redes neuronales se remonta a hace varias décadas (el descenso de gradiente estocástico de 1952, el perceptrón de 1958, la retropropagación de 1986, las redes convolucionales profundas de 1995), su verdadero auge se debe a la maduración simultánea de tres fuerzas impulsoras:

\begin{figure}[H]
\centering
\includegraphics[width=0.85\textwidth]{slides-images/slide-019.jpg}
\caption{La historia de las redes neuronales y las tres fuerzas impulsoras del auge del aprendizaje profundo\protect\footnotemark}
\end{figure}
\footnotetext{Diapositivas, página 19.}

\begin{table}[H]
\centering
\begin{tabular}{lp{10cm}}
\toprule
\textbf{Fuerza impulsora} & \textbf{Manifestación concreta} \\
\midrule
Big Data & La aparición de conjuntos de datos a gran escala (como ImageNet, Wikipedia) y la fuerte caída del costo de recolección y almacenamiento de datos \\
Hardware & La capacidad de cómputo masivamente paralela de las GPU hizo posible entrenar redes profundas \\
Software & Los marcos de código abierto (TensorFlow, PyTorch, Keras) simplificaron la construcción y el entrenamiento de modelos de aprendizaje profundo \\
\bottomrule
\end{tabular}
\caption{Las tres fuerzas impulsoras del auge del aprendizaje profundo}
\end{table}

\begin{knowledgebox}{Perspectiva histórica: estas técnicas no son nuevas}
Amini enfatiza en particular que casi todas las técnicas presentadas en esta clase (el perceptrón, la retropropagación, el descenso de gradiente) se inventaron hace \textbf{varias décadas}. Lo verdaderamente nuevo es la capacidad asombrosa que muestran cuando confluyen big data, un poder de cómputo considerable y buenas herramientas. Comprender estos fundamentos es esencial para seguir haciendo avanzar el campo.
\end{knowledgebox}

\subsection{Resumen de la sección}

La ventaja central del aprendizaje profundo frente al aprendizaje automático tradicional consiste en aprender automáticamente representaciones de características, sin necesidad de ingeniería manual de características. Aunque la teoría de base tiene ya varias décadas de historia, la triple fuerza impulsora de big data, hardware de GPU y software de código abierto ha hecho que el aprendizaje profundo viva un desarrollo explosivo en los últimos años.

